\documentclass[11pt,a4paper]{article}

\usepackage[T1]{fontenc}
\usepackage[utf8]{inputenc}
\usepackage[polish]{babel}
\usepackage{amsmath,amssymb,mathtools}
\usepackage{geometry}
\usepackage{hyperref}

\geometry{margin=2.5cm}

\newcommand{\C}{\mathbb{C}}
\newcommand{\Z}{\mathbb{Z}}
\newcommand{\id}{\mathbb{I}}
\newcommand{\Eold}{E_{\mathrm{old}}}
\newcommand{\Enew}{E_{\mathrm{new}}}
\newcommand{\Cold}{\mathcal{C}_{\mathrm{old}}}
\newcommand{\Cnew}{\mathcal{C}_{\mathrm{new}}}
\newcommand{\w}{\omega}
\newcommand{\CZ}{\operatorname{CZ}}
\newcommand{\diag}{\operatorname{diag}}
\newcommand{\im}{\operatorname{im}}
\newcommand{\rank}{\operatorname{rank}}
\newcommand{\ket}[1]{\left|#1\right\rangle}
\newcommand{\bra}[1]{\left\langle#1\right|}
\newcommand{\braket}[1]{\left\langle#1\right\rangle}

\title{Zmiana bazy kodowania w qutrytowych stanach grafowych}
\author{}
\date{}

\begin{document}
\maketitle

\section{Przestrzeń lokalna i stare kodowanie}

W projekcie jeden qutryt logiczny
\[
    \mathcal{H}_L \simeq \C^3
\]
jest kodowany w dwóch kubitach fizycznych
\[
    \mathcal{H}_{\mathrm{phys}}
    =
    \C^2 \otimes \C^2
    \simeq
    \C^4 .
\]
Używamy porządku bazy fizycznej
\[
    \ket{00},\ket{01},\ket{10},\ket{11}.
\]
Stare, referencyjne kodowanie jest izometrią
\[
    \Eold:\C^3\longrightarrow \C^4
\]
daną przez
\[
    \Eold
    =
    \begin{pmatrix}
        1 & 0 & 0\\
        0 & 1 & 0\\
        0 & 0 & 1\\
        0 & 0 & 0
    \end{pmatrix},
    \qquad
    \Eold^\dagger \Eold = \id_3 .
\]
Oznacza to
\[
    \ket{0}_L\mapsto \ket{00},\qquad
    \ket{1}_L\mapsto \ket{01},\qquad
    \ket{2}_L\mapsto \ket{10}.
\]
Stara przestrzeń kodowa to
\[
    \Cold
    =
    \im(\Eold)
    =
    \operatorname{span}\{\ket{00},\ket{01},\ket{10}\},
\]
a niewykorzystany kierunek dopełniający to
\[
    \ket{n_{\mathrm{old}}}=\ket{11}.
\]
Projektor na starą przestrzeń kodową ma postać
\[
    P_{\mathrm{old}}
    =
    \Eold\Eold^\dagger
    =
    \ket{00}\bra{00}
    +
    \ket{01}\bra{01}
    +
    \ket{10}\bra{10}.
\]

\section{Ogólna zmiana kodowania}

Nowe kodowanie jest macierzą
\[
    \Enew\in\C^{4\times 3}
\]
spełniającą warunek izometrii
\[
    \Enew^\dagger \Enew = \id_3.
\]
Jej kolumny są nowymi codewordami:
\[
    \Enew
    =
    \begin{pmatrix}
        \vert & \vert & \vert\\
        \ket{e_0} & \ket{e_1} & \ket{e_2}\\
        \vert & \vert & \vert
    \end{pmatrix},
    \qquad
    \braket{e_i|e_j}=\delta_{ij}.
\]
Dla dowolnego qutrytu logicznego
\[
    \ket{\psi}_L
    =
    \alpha_0\ket{0}
    +
    \alpha_1\ket{1}
    +
    \alpha_2\ket{2}
\]
stan fizyczny w nowym kodowaniu wynosi
\[
    \ket{\psi}_{\Enew}
    =
    \Enew\ket{\psi}_L
    =
    \alpha_0\ket{e_0}
    +
    \alpha_1\ket{e_1}
    +
    \alpha_2\ket{e_2}.
\]
Nowa przestrzeń kodowa to
\[
    \Cnew=\im(\Enew),\qquad
    P_{\mathrm{new}}=\Enew\Enew^\dagger.
\]

\subsection{Unitarny operator zmiany kodowania}

Ponieważ zarówno \(\Eold\), jak i \(\Enew\) są izometriami z \(\C^3\) do
\(\C^4\), można rozszerzyć zmianę bazy do pełnej unitarnej macierzy
\[
    W_{\Enew}\in U(4)
\]
takiej, że
\[
    W_{\Enew}\Eold=\Enew.
\]
W kodzie jest używany następujący wybór. Bierzemy znormalizowany wektor
\[
    \ket{n_{\mathrm{new}}}\in\ker(\Enew^\dagger),
    \qquad
    \Enew^\dagger\ket{n_{\mathrm{new}}}=0,
    \qquad
    \braket{n_{\mathrm{new}}|n_{\mathrm{new}}}=1.
\]
Wtedy
\[
    \boxed{
    W_{\Enew}
    =
    \Enew\Eold^\dagger
    +
    \ket{n_{\mathrm{new}}}\bra{n_{\mathrm{old}}}
    }.
\]
Pierwszy składnik przenosi starą trójwymiarową przestrzeń kodową na nową,
a drugi składnik przenosi brakujący kierunek \(\ket{11}\) na brakujący
kierunek \(\ket{n_{\mathrm{new}}}\).

Rzeczywiście,
\[
    W_{\Enew}\Eold
    =
    \Enew\Eold^\dagger\Eold
    +
    \ket{n_{\mathrm{new}}}\bra{n_{\mathrm{old}}}\Eold
    =
    \Enew\id_3+0
    =
    \Enew.
\]
Ponadto \(W_{\Enew}\) jest unitarne, bo odwzorowuje ortonormalną bazę
\[
    \{\ket{00},\ket{01},\ket{10},\ket{11}\}
\]
na ortonormalną bazę
\[
    \{\ket{e_0},\ket{e_1},\ket{e_2},\ket{n_{\mathrm{new}}}\}.
\]
Algebraicznie:
\[
    W_{\Enew}^\dagger W_{\Enew}
    =
    \Eold\Enew^\dagger\Enew\Eold^\dagger
    +
    \ket{n_{\mathrm{old}}}\bra{n_{\mathrm{new}}}
    \ket{n_{\mathrm{new}}}\bra{n_{\mathrm{old}}}
    =
    \Eold\Eold^\dagger
    +
    \ket{11}\bra{11}
    =
    \id_4.
\]

\section{Qutrytowy stan grafowy}

Dla grafu \(G=(V,E)\), \(|V|=n\), qutrytowy stan grafowy w przestrzeni
logicznej można zapisać jako
\[
    \ket{G}
    =
    \prod_{\{a,b\}\in E}\CZ_{ab}^{(3)}
    \ket{+}_3^{\otimes n},
\]
gdzie
\[
    \ket{+}_3
    =
    \frac{1}{\sqrt{3}}
    \sum_{j=0}^{2}\ket{j},
    \qquad
    \w=e^{2\pi i/3},
\]
oraz
\[
    \CZ^{(3)}\ket{i,j}
    =
    \w^{ij}\ket{i,j},
    \qquad
    i,j\in\Z_3.
\]
Jeżeli \(A_{ab}\) oznacza krotność krawędzi \(\{a,b\}\) modulo \(3\), to
równoważny zapis amplitudowy jest następujący:
\[
    \boxed{
    \ket{G_A}
    =
    \frac{1}{\sqrt{3^n}}
    \sum_{x\in\Z_3^n}
    \w^{q_A(x)}
    \ket{x_1,\ldots,x_n}
    },
    \qquad
    q_A(x)=\sum_{1\le a<b\le n} A_{ab}x_ax_b.
\]
To dobrze opisuje również przypadek grafów z powtórzonymi krawędziami: dwie
identyczne krawędzie dają dwa razy zastosowaną bramkę \(\CZ^{(3)}\), czyli
fazę \(\w^{2x_ax_b}\).

Po zakodowaniu każdego qutrytu przez tę samą izometrię \(E\) otrzymujemy
\[
    \ket{G_A;E}
    =
    E^{\otimes n}\ket{G_A}.
\]
W szczególności
\[
    \ket{G_A;\Enew}
    =
    \Enew^{\otimes n}\ket{G_A}
    =
    (W_{\Enew}\Eold)^{\otimes n}\ket{G_A}
    =
    W_{\Enew}^{\otimes n}\ket{G_A;\Eold}.
\]
To jest główna matematyczna treść zmiany bazy w benchmarku: idealny stan
fizyczny w nowym kodowaniu jest starym fizycznym stanem grafowym, na który na
każdym zakodowanym qutrycie działa lokalny operator \(W_{\Enew}\).

\section{Dwie równoważne strategie obwodowe}

Niech \(\widetilde{F}_3\) i \(\widetilde{\CZ}_3\) oznaczają fizyczne wersje
bramek qutrytowych w starym kodowaniu. Strategia \texttt{append\_w} buduje
najpierw standardowy obwód stanu grafowego w starym kodowaniu:
\[
    \ket{G_A;\Eold}
    =
    \left(\prod_{\{a,b\}\in E}\widetilde{\CZ}_{3,ab}\right)
    \widetilde{F}_3^{\otimes n}
    \ket{00}^{\otimes n},
\]
a następnie dodaje
\[
    W_{\Enew}^{\otimes n}
\]
na końcu.

Druga strategia, \texttt{prepared\_w\_then\_conjugated\_entanglers}, przesuwa
zmianę bazy do środka konstrukcji. Lokalny stan początkowy jest przygotowywany
jako
\[
    W_{\Enew}\ket{+}_3
    =
    W_{\Enew}\widetilde{F}_3\ket{00},
\]
a każdy entangler na krawędzi jest sprzężony:
\[
    \boxed{
    \widetilde{\CZ}_{\Enew}
    =
    (W_{\Enew}\otimes W_{\Enew})
    \widetilde{\CZ}_{3}
    (W_{\Enew}^\dagger\otimes W_{\Enew}^\dagger)
    }.
\]
Wtedy
\[
    \prod_{\{a,b\}\in E}
    \widetilde{\CZ}_{\Enew,ab}
    \left(W_{\Enew}\ket{+}_3\right)^{\otimes n}
    =
    W_{\Enew}^{\otimes n}
    \prod_{\{a,b\}\in E}
    \widetilde{\CZ}_{3,ab}
    \ket{+}_3^{\otimes n}.
\]
Idealnie obie strategie dają ten sam stan. Różnią się jednak rozkładem na
bramki po transpilacji, dlatego benchmark porównuje głębokość, rozmiar i liczbę
bramek dwukubitowych.

\section{Bazy monomialne}

\subsection{Definicja klasy \texttt{monomial\_full}}

W \texttt{encoding\_search\_v2} klasa \texttt{monomial\_full} wybiera trzy z
czterech stanów bazy fizycznej. Niech
\[
    S=\{s_0,s_1,s_2\}\subset\{0,1,2,3\},
    \qquad |S|=3,
\]
gdzie indeksy oznaczają
\[
    0\leftrightarrow \ket{00},\quad
    1\leftrightarrow \ket{01},\quad
    2\leftrightarrow \ket{10},\quad
    3\leftrightarrow \ket{11}.
\]
Niech
\[
    B_S=
    \begin{pmatrix}
        \vert & \vert & \vert\\
        \ket{s_0} & \ket{s_1} & \ket{s_2}\\
        \vert & \vert & \vert
    \end{pmatrix}
    \in\C^{4\times 3}
\]
będzie macierzą wyboru supportu. Następnie wybieramy permutację
\[
    \pi\in S_3
\]
oraz fazy
\[
    \lambda_0,\lambda_1,\lambda_2\in\{1,\w,\w^2\}.
\]
Macierz permutacji w konwencji kodu spełnia
\[
    (P_\pi)_{r,\pi_r}=1,
\]
a macierz faz to
\[
    D_\lambda=\diag(\lambda_0,\lambda_1,\lambda_2).
\]
Kandydat monomialny ma postać
\[
    \boxed{
    E_{S,\pi,\lambda}
    =
    B_S D_\lambda P_\pi
    }.
\]
Jest to izometria, bo
\[
    E_{S,\pi,\lambda}^\dagger E_{S,\pi,\lambda}
    =
    P_\pi^\dagger D_\lambda^\dagger B_S^\dagger B_S D_\lambda P_\pi
    =
    P_\pi^\dagger D_\lambda^\dagger D_\lambda P_\pi
    =
    \id_3.
\]

Liczba takich kandydatów wynosi
\[
    \binom{4}{3}\cdot 3!\cdot 3^3
    =
    4\cdot 6\cdot 27
    =
    648.
\]

\subsection{Interpretacja codewordów}

Każdy kandydat monomialny przypisuje każdemu symbolowi logicznemu dokładnie
jeden stan bazy fizycznej, z lokalną fazą. Można więc zapisać
\[
    E_{S,\pi,\lambda}\ket{j}
    =
    \alpha_j\ket{m(j)},
    \qquad
    \alpha_j\in\{1,\w,\w^2\},
\]
gdzie
\[
    m:\{0,1,2\}\hookrightarrow \{00,01,10,11\}
\]
jest injekcją wyznaczoną przez support i permutację.

Dla stanu grafowego daje to jawny wzór
\[
    \boxed{
    \ket{G_A;E_{S,\pi,\lambda}}
    =
    \frac{1}{\sqrt{3^n}}
    \sum_{x\in\Z_3^n}
    \w^{q_A(x)}
    \left(\prod_{v=1}^n \alpha_{x_v}\right)
    \ket{m(x_1),\ldots,m(x_n)}
    }.
\]
Zmiana monomialna nie miesza symboli w superpozycje fizyczne. Ona tylko:
\[
    \text{wybiera support }3\text{ z }4,\quad
    \text{permutuje etykiety logiczne},\quad
    \text{dokłada fazy lokalne}.
\]
Dlatego pojedyncze codewordy monomialne są zawsze stanami produktowymi bazy
obliczeniowej dwóch kubitów i mają zerowe splątanie wewnątrz jednego
zakodowanego qutrytu.

\subsection{Stara przestrzeń kodowa jako szczególny przypadek}

Jeżeli
\[
    S=\{0,1,2\},
\]
to
\[
    B_S=\Eold
\]
i dostajemy starszą klasę monomialną ograniczoną do starej code space:
\[
    E_{\mathrm{old\;monomial}}
    =
    \Eold D_\lambda P_\pi.
\]
Wtedy nowa przestrzeń kodowa jest równa starej:
\[
    \im(E_{\mathrm{old\;monomial}})=\Cold.
\]
Zmienia się tylko baza logiczna wewnątrz tej samej trójwymiarowej przestrzeni.
Dla takiego przypadku można wybrać
\[
    W
    =
    \Eold D_\lambda P_\pi \Eold^\dagger
    +
    \ket{11}\bra{11},
\]
czyli macierz blokową
\[
    W
    =
    \begin{pmatrix}
        D_\lambda P_\pi & 0\\
        0 & 1
    \end{pmatrix}
\]
w rozkładzie
\[
    \C^4=\Cold\oplus\operatorname{span}\{\ket{11}\}.
\]

\subsection{Gdy support zawiera \(\ket{11}\)}

Jeżeli support \(S\) zawiera indeks \(3\), czyli stan \(\ket{11}\), to nowa
przestrzeń kodowa nie jest już równa starej. Na przykład dla
\[
    S=\{0,1,3\}
\]
używanymi stanami fizycznymi są
\[
    \ket{00},\ket{01},\ket{11},
\]
a stan \(\ket{10}\) trafia do dopełnienia. Nadal każdy codeword jest prostym
stanem bazy obliczeniowej, ale sama trójwymiarowa przestrzeń kodowa jest innym
podzbiorem \(\C^4\).

Wtedy \(W\) jest pełną czterowymiarową permutacją z fazami: przenosi
\[
    \ket{00},\ket{01},\ket{10}
\]
na trzy wybrane stany supportu, a \(\ket{11}\) na jedyny brakujący stan bazy
fizycznej, z fazą wynikającą z wyboru wektora dopełniającego.

\section{Bazy produktowe}

\subsection{Definicja klasy \texttt{product}}

W klasie \texttt{product} wybieramy dwie lokalne unitarne macierze
\[
    U,V\in U(2)
\]
i definiujemy
\[
    \boxed{
    E_{U,V}
    =
    (U\otimes V)\Eold
    }.
\]
W wersji dyskretnej \texttt{encoding\_search\_v2} bierze
\[
    U,V\in
    \{I,X,SX,SX^\dagger,H,HS,SH,Z\}.
\]
W opcjonalnym gridzie używana jest parametryzacja
\[
    U(\alpha,\beta,\gamma)
    =
    R_z(\alpha)R_x(\beta)R_z(\gamma),
\]
gdzie
\[
    R_x(\beta)
    =
    \begin{pmatrix}
        \cos(\beta/2) & -i\sin(\beta/2)\\
        -i\sin(\beta/2) & \cos(\beta/2)
    \end{pmatrix},
\]
\[
    R_z(\alpha)
    =
    \begin{pmatrix}
        e^{-i\alpha/2} & 0\\
        0 & e^{i\alpha/2}
    \end{pmatrix}.
\]

\subsection{Izometria i dokładne \(W\)}

Ponieważ \(U\otimes V\) jest unitarną macierzą \(4\times 4\), mamy
\[
    E_{U,V}^\dagger E_{U,V}
    =
    \Eold^\dagger (U^\dagger\otimes V^\dagger)(U\otimes V)\Eold
    =
    \Eold^\dagger\Eold
    =
    \id_3.
\]
Dla tej klasy naturalnym operatorem zmiany kodowania jest po prostu
\[
    \boxed{
    W_{U,V}=U\otimes V
    }.
\]
Rzeczywiście,
\[
    W_{U,V}\Eold=(U\otimes V)\Eold=E_{U,V}.
\]
Wektor dopełniający jest wtedy jawny:
\[
    \ket{n_{U,V}}
    =
    (U\otimes V)\ket{11}
    =
    U\ket{1}\otimes V\ket{1}.
\]

\subsection{Codewordy produktowe}

Dla starej bazy logicznej mamy
\[
    \ket{0}_L\mapsto \ket{00},
    \qquad
    \ket{1}_L\mapsto \ket{01},
    \qquad
    \ket{2}_L\mapsto \ket{10}.
\]
Po bazie produktowej:
\[
    \ket{0}_L
    \mapsto
    U\ket{0}\otimes V\ket{0},
\]
\[
    \ket{1}_L
    \mapsto
    U\ket{0}\otimes V\ket{1},
\]
\[
    \ket{2}_L
    \mapsto
    U\ket{1}\otimes V\ket{0}.
\]
Każdy pojedynczy codeword jest stanem separowalnym dwóch kubitów. Dlatego
wewnętrzne splątanie codewordów jest równe zero:
\[
    S\!\left(\operatorname{Tr}_2\ket{e_j}\bra{e_j}\right)=0,
    \qquad j=0,1,2.
\]
To odróżnia \texttt{product} od ogólnych izometrii \(4\times 3\), ale nie od
monomiali, które też mają niesplątane pojedyncze codewordy. Różnica polega na
tym, że produktowa baza może tworzyć superpozycje stanów obliczeniowych, ale
wyłącznie przez lokalną rotację dwóch kubitów.

\subsection{Stan grafowy w bazie produktowej}

Dla stanu grafowego
\[
    \ket{G_A;E_{U,V}}
    =
    E_{U,V}^{\otimes n}\ket{G_A}
\]
dostajemy
\[
    \boxed{
    \ket{G_A;E_{U,V}}
    =
    (U\otimes V)^{\otimes n}
    \ket{G_A;\Eold}
    }.
\]
Po rozwinięciu w amplitudach:
\[
    \ket{G_A;E_{U,V}}
    =
    \frac{1}{\sqrt{3^n}}
    \sum_{x\in\Z_3^n}
    \w^{q_A(x)}
    \bigotimes_{v=1}^{n}
    (U\otimes V)\ket{b_{x_v}},
\]
gdzie
\[
    \ket{b_0}=\ket{00},\qquad
    \ket{b_1}=\ket{01},\qquad
    \ket{b_2}=\ket{10}.
\]
W przeciwieństwie do monomiali, pojedynczy składnik sumy może być superpozycją
w fizycznej bazie \(\ket{00},\ket{01},\ket{10},\ket{11}\). Cały stan nadal
leży jednak w przestrzeni
\[
    \left((U\otimes V)\Cold\right)^{\otimes n}.
\]

\subsection{Overlap ze starą code space}

Metryka używana w benchmarku ma postać
\[
    \operatorname{overlap}(E,\Cold)
    =
    \frac{1}{3}\|P_{\mathrm{old}}E\|_F^2.
\]
Dla \(E_{U,V}=(U\otimes V)\Eold\) można ją zapisać jako
\[
    \operatorname{overlap}(E_{U,V},\Cold)
    =
    1-\frac{1}{3}
    \sum_{j=0}^{2}
    \left|
    \bra{11}(U\otimes V)\ket{b_j}
    \right|^2.
\]
Jeżeli
\[
    U=(u_{ab})_{a,b=0}^{1},
    \qquad
    V=(v_{ab})_{a,b=0}^{1},
\]
to
\[
    \bra{11}(U\otimes V)\ket{00}=u_{10}v_{10},
\]
\[
    \bra{11}(U\otimes V)\ket{01}=u_{10}v_{11},
\]
\[
    \bra{11}(U\otimes V)\ket{10}=u_{11}v_{10}.
\]
Stąd
\[
    \operatorname{overlap}(E_{U,V},\Cold)
    =
    1-\frac{1}{3}
    \left(
        |u_{10}v_{10}|^2
        +
        |u_{10}v_{11}|^2
        +
        |u_{11}v_{10}|^2
    \right).
\]
Dla \(U=V=I\) overlap wynosi \(1\). Dla lokalnych rotacji mieszających
\(\ket{0}\) i \(\ket{1}\) część amplitudy przechodzi do kierunku \(\ket{11}\),
więc overlap ze starą code space maleje.

\section{Porównanie: \texttt{monomial} kontra \texttt{product}}

\[
\begin{array}{c|c|c}
    \text{cecha} & \texttt{monomial} & \texttt{product}\\
    \hline
    \text{postać} &
    B_SD_\lambda P_\pi &
    (U\otimes V)\Eold\\
    \text{codewordy} &
    \text{stany bazy obliczeniowej z fazami} &
    \text{lokalne stany produktowe}\\
    \text{superpozycje lokalne} &
    \text{nie} &
    \text{tak, ale separowalne}\\
    \text{splątanie codewordu} &
    0 &
    0\\
    \text{może użyć }\ket{11}\text{ jako codewordu} &
    \text{tak, dyskretnie} &
    \text{tak, przez lokalną rotację}\\
    \text{naturalne }W &
    \text{permutacja z fazami w }\C^4 &
    U\otimes V\\
    \text{koszt splątujący samego }W &
    \text{zależy od permutacji} &
    0
\end{array}
\]

Najważniejsza różnica praktyczna jest taka:
\[
    \texttt{monomial}
    \quad
    \text{testuje dyskretne relabelingi i fazy symboli qutrytowych,}
\]
podczas gdy
\[
    \texttt{product}
    \quad
    \text{testuje lokalne rotacje dwóch kubitów tworzących jeden qutryt.}
\]
Obie klasy zachowują prostotę pojedynczych codewordów: żaden codeword sam w
sobie nie jest splątany między dwoma kubitami kodującymi qutryt. Różnią się
tym, czy zmiana jest dyskretną zmianą supportu i faz, czy ciągłą/lokalną
rotacją fizycznych kubitów.

\section{Kryterium trywialności w pipeline}

Pipeline zachowuje referencję
\[
    \texttt{baseline}/\texttt{E\_old}
\]
i pomija kandydatów równoważnych bazie referencyjnej. Matematycznie są to
przypadki, w których
\[
    \Enew=\Eold
\]
w tolerancji numerycznej, albo
\[
    \Enew=e^{i\phi}\Eold
\]
do jednej globalnej fazy, albo zbudowany operator
\[
    W_{\Enew}
\]
jest identycznością do tolerancji. Takie przypadki nie zmieniają fizycznego
problemu benchmarkowego, więc nie warto ponownie ich transpilować.

\section{Podsumowanie}

Cała zmiana bazy w badaniu stanów grafowych sprowadza się do wyboru izometrii
\[
    \Enew:\C^3\to\C^4
\]
i unitarnego rozszerzenia
\[
    W_{\Enew}\in U(4)
    \quad\text{takiego, że}\quad
    W_{\Enew}\Eold=\Enew.
\]
Dla \(n\)-qutrytowego stanu grafowego:
\[
    \boxed{
    \ket{G;\Enew}
    =
    W_{\Enew}^{\otimes n}\ket{G;\Eold}
    }.
\]
Dla bazy monomialnej \(W_{\Enew}\) jest zasadniczo permutacją z fazami w
czterowymiarowej bazie fizycznej. Dla bazy produktowej
\[
    W_{\Enew}=U\otimes V,
\]
czyli zmiana kodowania jest lokalna względem dwóch kubitów pojedynczego
qutrytu. Obie klasy są więc zrozumiałymi, kontrolowanymi podrodzinami ogólnych
izometrii \(4\times 3\), ale sondują inne mechanizmy upraszczania obwodu po
transpilacji.

\end{document}
